\section{Мысленные эксперименты}

\subsection{Писец}
Писец пишет по одной цифре каждый планковский миг --- быстрее не меняется ничто во
Вселенной. С Большого взрыва он успел бы вывести около \(10^{60}\) цифр. У уровня \(b\) их
\(\bvalue\). Он не закончил и ничтожной доли; конца \(b\) он не достигнет никогда --- не
говоря уже о \(c\). А \(b\) --- самый маленький из наших великанов. Если даже время не
доходит до конца \(b\), что значит сказать, что у \(b\) «столько-то» цифр?

\subsection{Вавилонская библиотека}
Библиотека Борхеса хранит все книги по 410 страниц --- вообразимо огромная, но всё же
конечная. Рядом с \(d\) она смехотворно мала. Чтобы расставить по полкам цифры \(d\), нужна
библиотека, каталог которой нельзя записать в этой Вселенной. Борхес придумал самую большую
библиотеку, какую человек в силах вынести; \(d\) требует той, что вынести нельзя.

\subsection{Тот, кто мог бы его увидеть}
Вообразите ум, способный удержать крипто кому целиком --- все цифры разом. Была бы для него
крипто кома одним числом, как семёрка для нас? Или «одно число» на таком масштабе перестаёт
что-либо значить? Быть может, по-настоящему увидеть число --- значит сделать его маленьким.
А понимание --- не другое ли это слово для «уменьшить»?
